\section{C. Sections over a base and the cusp extension}
\label{part:C}

We use the following notation throughout this part. A smooth pair over a scheme
$S$ is a smooth proper curve $\bar U\to S$ with geometrically connected fibres,
together with an effective relative divisor $D\subset\bar U$ finite and
\emph{\'etale} over $S$; we put $U=\bar U\setminus D$. Thus collisions of
marked points are excluded from the term ``smooth pair''. The type of a fibre
is $(g,r)$, with $r=\deg(D/S)$. We fix compatible geometric base points when
writing homomorphisms of fundamental groups. Changing a path to a base point
conjugates a displayed homomorphism by an element of the indicated geometric
kernel. We write $G=\pi_1(S)$, $E=\pi_1(U)$ and
$\Pi=\pi_1(U_{\bar s})$. These are respectively $E_S$, $E_U$ and $\pi$ in
\S3 of the manuscript, scans 64 and 79. In residue characteristic $p>0$, a
superscript $(p')$ on a geometric kernel means its maximal pro-prime-to-$p$
quotient; we retain the full arithmetic quotient. We write
$T_{p'}=\prod_{\ell\ne p}\mathbb Z_\ell(1)$, and write $T$ when $p=0$ and
all primes occur. All cohomology of profinite groups below is continuous.

\begin{result}{C1}{Proved}{The Hartogs extension of a rational section}
\textup{[F11; \S3, manuscript pp.\ 32, 34--35, 41; scans 65, 70--72, 83.]}
Let $S$ be a smooth integral scheme of finite type over a field of
characteristic zero. Let $U\to S$ be a smooth pair of constant type different
from $(0,0)$ and $(0,1)$. For a closed subset $Z\subset S$ of codimension at
least two, restriction is a bijection
\[
 \operatorname{Sect}_S(U)\longrightarrow
 \operatorname{Sect}_{S\setminus Z}(U\times_S(S\setminus Z)).
\]
The same assertion holds for a finite tower of such relative curves.
\end{result}

\begin{proof}
Uniqueness follows from separatedness: two sections agree on the schematically
dense open set $S\setminus Z$, hence their equalizer, the inverse image of the
diagonal, is all of the reduced scheme $S$.

First suppose $g\geq1$. We give the resolution argument behind the manuscript's
reference to Weil. The question is local for the \'etale topology of $S$.
A smooth proper curve has a section after an \'etale localization, and then
a divisor of sufficiently large relative degree gives a projective embedding.
We can consequently resolve the rational section $S\dashrightarrow\bar U$
by principalizing its base ideal. Bierstone--Milman, Theorem 1.10, supplies
a finite sequence of blowups with smooth centres, disjoint from the locus
where the rational map is defined, on which the ideal becomes invertible
\cite{CF_BierstoneMilman1997}. The resolved map is an $S$-morphism.

Consider the last blowup $b:S_m\to S_{m-1}$. Its geometric fibres are points
or projective spaces. Every morphism from $\mathbb P^1$ to a smooth proper
curve of genus at least one is constant: a nonconstant morphism would be
finite, and the Riemann--Hurwitz equality
$-2=d(2g-2)+\deg R$ is impossible. Since any two points of a projective
space lie on a projective line, the resolved map is constant on every
geometric fibre of $b$. Let $\Gamma$ be its scheme-theoretic graph image in
$S_{m-1}\times_S\bar U$. Projection $\Gamma\to S_{m-1}$ is proper,
birational, and has finite fibres, hence is finite. The target is normal,
so this finite birational morphism is an isomorphism: on an affine open,
its coordinate ring is an integral subalgebra of the common fraction field
and is contained in the integrally closed coordinate ring of the target.
The map therefore descends across $b$. Descending successively gives a
section $\bar f:S\to\bar U$. The inverse image $\bar f^{-1}D$ is the zero
scheme of a nonzero section of an invertible sheaf. If nonempty, it has a
codimension-one component, by the principal ideal theorem in the regular
local rings of $S$. Its support lies in $Z$, so it is empty.

For $g=0$, the hypothesis gives $r\geq2$. After an \'etale localization a
boundary section is defined. A genus-zero smooth proper curve with this
section is locally a projective-line bundle; locally on $S$, its complement
is $\mathbb A^1_S$. The given section is a regular function on $S\setminus Z$.
The normality equality
$\mathcal O_S=j_*\mathcal O_{S\setminus Z}$ extends that function to $S$.
For each remaining boundary section, its intersection with the extended
section is again an effective Cartier divisor supported in $Z$, and is empty.
The unique local extensions descend.

For a tower, first extend the section of the bottom curve. Pull the next
curve back along that extended section, apply the one-curve assertion, and
continue. These are actual Cartesian base changes, and the restriction of
each pullback to $S\setminus Z$ is the pullback already carrying the given
section. This proves the assertion without assuming that the entire tower
embeds in a semiabelian variety.

The exclusion of $(0,0)$ matters: the section
$(x,y)\mapsto[x:y]$ of $\mathbb P^1\times(\mathbb A^2\setminus\{0\})$
does not extend at the origin. Its restrictions to the two coordinate axes
have different constant limits. The genus-zero proof also proves Hartogs
extension for an affine-line bundle; that observation does not give the
inertia criterion below, which fails for type $(0,1)$.
\end{proof}

\begin{result}{C2}{Proved}{The trait criterion, including the two nonhyperbolic types}
\textup{[F12; \S3, manuscript pp.\ 36--38; scans 74--78.]}
Let $R$ be a discrete valuation ring, $K$ its fraction field, and let
$U=\bar U\setminus D$ be a smooth pair over $R$ of type other than
$(0,0),(0,1)$. Let $f_K\in U(K)$. In characteristic zero for the residue
field use all geometric covers; in residue characteristic $p>0$ use the
geometrically pro-prime-to-$p$ fundamental groups. The following are
equivalent:
\begin{enumerate}
\item $f_K$ extends to an $R$-section of $U$;
\item the continuous lift $G_K\to E_U^{(p')}$ induced by $f_K$ factors
through $\pi_1(\operatorname{Spec}R)$;
\item this lift is trivial on an inertia group at the closed point.
\end{enumerate}
Here ``factors'' means that the induced map on the quotient is a section
of $E_U^{(p')}\to\pi_1(R)$. For nonhenselian $R$, inertia is taken in a
chosen decomposition group, and its closed normal closure is the kernel
of $G_K\to\pi_1(R)$. The statement extends to a tower of smooth pairs
whose curve factors have these types.
\end{result}

\begin{proof}
We can pass to the strict henselization to test the last condition. Extension
of the rational point can also be tested there: properness first extends it
uniquely to $\bar U$ over $R$, and whether its closed-point image lies in $D$
is unchanged by faithfully flat base change. Thus the descent in this step
is descent of this specified proper extension, not descent of an arbitrarily
chosen section.

If its closed-point image is in $U$, functoriality gives the required
factorization. Otherwise it is a cusp $q$. C3 below identifies the inertia
map with $\kappa_q\circ[\mu]$, where $\mu>0$ and $\kappa_q$ is injective.
Multiplication by a nonzero integer is injective and nonzero on
$T_{p'}$. The inertia map is consequently nontrivial. This proves the
one-curve assertion. The peripheral injection used here follows from the
prime-to-$p$ presentation of a punctured smooth curve in SGA~1,
Expos\'e XIII, Corollary 2.12, and its peripheral-subgroup form recorded in
Stix, \S2 and Corollary 12 \cite{CF_SGA1,CF_Stix2008}. For a single cusp
the group can be nonabelian with the cusp zero in its abelianization; C3
does not abelianize it.

For reference, the needed peripheral injection can be checked in every
pro-$\ell$ quotient directly. The presentation of an affine curve is the
free pro-$\ell$ group on $a_1,b_1,\ldots,a_g,b_g,c_1,\ldots,c_{r-1}$,
with the last cusp the inverse of
$\prod_j[a_j,b_j]\prod_i c_i$. For $r\geq2$ each peripheral word has
nonzero abelianization after choosing that cusp among the free generators.
For $r=1,g\geq1$ its degree-two class is
$\sum_j[a_j,b_j]\ne0$ in the free degree-two Lie module. It is therefore
nontrivial and has infinite order in this torsion-free free pro-$\ell$
group. The homomorphism $\mathbb Z_\ell(1)$ is injective. Taking all
$\ell\ne p$ detects every coordinate of $T_{p'}$ and proves the required
injection, including the one-cusp case.

For type $(0,2)$ the cusp is the whole geometric group $T_{p'}$, so the
same proof applies. For $(1,0)$ properness already gives extension into
$U=\bar U$, and smooth proper specialization gives trivial inertia in the
indicated quotient. No centre-free hypothesis is needed for either case.
For $(0,1)$, $R=k[[t]]$ and $f_K=t^{-1}\in\mathbb A^1(K)$ give a
counterexample: the prime-to-$p$ geometric group of the affine line is
trivial, whereas the point has a pole.

For a tower, project the lift to the first curve. The one-curve criterion
extends that projected point. Pull back the next curve to $\operatorname{Spec}R$
along this extended section. The fundamental-group square is Cartesian:
both its horizontal maps induce the same isomorphism on geometric kernels,
and the right-hand quotient is $G_K\to\pi_1(R)$; its exact sequences give
the Cartesian assertion by the elementary kernel-and-quotient argument.
For the exceptional centre-bearing types over a trait the same exactness
follows from the smooth proper or tame homotopy sequence with section,
SGA~1 XIII, Proposition 4.3 and Examples 4.4. The induced lift to this
pullback therefore kills inertia. Apply the one-curve criterion again and
iterate. This also proves exactly which comparison maps support the
manuscript's d\'evissage.
\end{proof}

\begin{result}{C3}{Proved}{The nonabelian intersection formula and the five conditions}
\textup{[F12; \S3, manuscript pp.\ 37--39 and 44; scans 75--79, 90.]}
Let $R$ be henselian and use the notation of C2. Suppose that the proper
extension of $f_K$ meets a rational cusp $q$ in the closed fibre. If $z$ is
a local equation of $q$ and $t$ a uniformizer of $R$, write
\[
 f_K^*z=t^\mu u,\qquad u\in R^*,\qquad \mu>0.
\]
With positive Kummer inertia defined by
$\tau\mapsto\tau(z^{1/n})/z^{1/n}$, the induced continuous homomorphism on
inertia is, up to the geometric conjugation arising from the chosen path,
\[
 I_K\longrightarrow I_K/P_K=T_{p'}
       \xrightarrow{[\mu]}T_{p'}\xrightarrow{\kappa_q}\Pi^{(p')}.
 \tag{C.1}
\]
For $p=0$, $P_K=1$. The formula holds with one cusp. With at least two
cusps it induces on the toric part of homology
\[
 T_\ell\longrightarrow T_\ell^{D(\bar k)}/T_\ell,
 \qquad a\longmapsto\mu a e_q \pmod{T_\ell\cdot(1,\ldots,1)}.
 \tag{C.2}
\]
In equicharacteristic zero, the five conditions a)--e) of scan 90 are
equivalent for the lift $\lambda:G_K\to E_U$: injectivity of $\lambda$;
the positive-multiplicity formula (C.1); nontrivial inertia; failure of
descent to $\pi_1(R)$; and failure of $f_K$ to extend into $U$. In that
case its image has index $\mu$ in the cusp normalizer. For $p>0$, replace
the domain in the injectivity assertion by $G_K/P_K$, and replace the
index by
$\mu_{p'}=\prod_{\ell\ne p}\ell^{v_\ell(\mu)}$.
\end{result}

\begin{proof}
The local ring of $\bar U$ at $q(s)$ is regular of dimension two and $z=0$
is a regular divisor. Over its strict henselization the prime-to-$p$ tame
covers of its complement are dominated by $z^{1/n}$, with $(n,p)=1$.
This is precisely SGA~1 XIII, Proposition 5.1 together with its
normal-crossings formulation, Proposition 5.2: adjoining prime-to-residue-
characteristic roots of regular parameters makes the cover \'etale
\cite{CF_SGA1}. These statements apply here to the one divisor $z=0$.
In the strict henselization of $R$, the unit $u$ has an $n$th root, by
Hensel's lemma. Compatible such roots exist over the inverse system
(its nonempty finite levels have surjective transition maps). Pulling
back $z^{1/n}$ along $f_K$ gives $t^{\mu/n}u^{1/n}$; its inertia character
is $\mu$ times the character of $t^{1/n}$. Passing through the local
fundamental group and then into the global group proves (C.1). Every
stage is a homomorphism of fundamental groups; no abelianization of the
target has occurred. Formula (C.2) follows by applying the specified
abelianization map and the single relation that the sum of peripheral
classes is zero. For one cusp this latter class is zero; the injection
$\kappa_q$ before abelianization remains injective.

The equality of intersection multiplicities in scans 77--78 is also
literal. If $A=\mathcal O_{\bar U,q(s)}$ and $I_f,I_q$ are the ideals of
the two sections, then
\[
 R/q^*I_f\simeq A/(I_f+I_q)\simeq R/f^*I_q.
\]
The lengths of these Artinian modules are equal to $\mu$.

For scan 90, the composite $G_K\xrightarrow\lambda E_U\to G_k$ is the
residue quotient. Hence $\ker\lambda\subset I_K$. In characteristic zero,
(C.1) is injective on $I_K$ whenever the section meets the boundary; this
proves injectivity of $\lambda$. If it does not meet the boundary, it
kills $I_K\ne1$ and is not injective. C2 proves the other equivalences.
The image normalizes $\mu L_q$. In a hyperbolic curve the normalizer of
any open subgroup of a peripheral group is the cusp normalizer, by the
peripheral self-normalization theorem; in type $(0,2)$ the entire
geometric group is peripheral and the same conclusion is immediate.
Alternatively, the map from the local punctured tube constructed above
lands directly in this normalizer. Its intersection with $L_q$ is
$\mu L_q$ and its image in $G_k$ is all of $G_k$. The coset index is
therefore $[T:\mu T]=\mu$. With prime-to-$p$ kernels the same proof
applies after quotienting $P_K$, and gives $[T_{p'}:\mu T_{p'}]=\mu_{p'}$.
Full $G_K$ need not inject in this version: wild inertia has been killed.

For the proposed test $\mathbb P^1_R\setminus\{0,\infty,q\}$ with
$q=t^\mu u$, the sections $0$ and $q$ collide. Thus its marked boundary
is not \'etale and it is not a smooth pair over $R$. The local Kummer
calculation $z\mapsto t^\mu u$ still tests the factor $\mu$ at the cusp
$0$. A test within the precise hypotheses is the fixed smooth pair
$\mathbb P^1_R\setminus\{0,1,\infty\}$ and the point $z=t^\mu u$.
Both tests have trivial normal line at $0$, so neither distinguishes
$c_1(N)$ from $-c_1(N)$. The global transition calculation in C6 fixes
that sign; assigning a nonzero Chern sign to a trivial line would not
test it.
\end{proof}

\begin{result}{C4}{Proved}{The exact sequence and the centre-free extension}
\textup{[F13; \S3, manuscript p.\ 39; scan 79; also scans 64, 73.]}
Let $S$ be a connected normal scheme of finite type over a field of
characteristic zero, and let $U\to S$ be a hyperbolic smooth pair. Then
\[
 1\longrightarrow\Pi\longrightarrow E\longrightarrow G\longrightarrow1
 \tag{C.3}
\]
is exact. Its formation commutes with the base changes to a normal
connected scheme used in this part, with the chosen geometric fibre
identified by a path. If $\rho:G\to\operatorname{Out}(\Pi)$ is its
continuous outer action, the continuous map
\[
 E\longrightarrow
 E_\rho=\{(a,g)\in\operatorname{Aut}(\Pi)\times G:
               [a]=\rho(g)\},\qquad
 e\longmapsto(\operatorname{conj}_e,\bar e)
 \tag{C.4}
\]
is an isomorphism. For a tower with factors also of types $(0,2)$ or
$(1,0)$, exactness holds over an \'etale $K(\pi,1)$ base, and at each
stage over a trait as in C2. It is false for arbitrary bases without
the hyperbolic or $\pi_2$ hypothesis.
\end{result}

\begin{proof}
Holzschuh--Schmidt--Stix, Theorem 5.6, applies to a geometric fibration
with geometrically connected fibres over a normal quasi-compact
quasi-separated scheme of characteristic zero. A smooth proper curve
with a relative normal-crossings boundary, in particular our finite
\'etale boundary, is such a fibration (their Definition 5.5). Their
quasifibration theorem gives the homotopy sequence. In the segment
\[
 \pi_2^{\mathrm{et}}(S)\xrightarrow\partial\Pi
                 \longrightarrow E\longrightarrow G\longrightarrow1,
\]
the image of $\partial$ is central in $\Pi$. This is the fibration
boundary property: the action of a fibre loop on a spherical lift is
trivial after the two boundary homotopies are composed; equivalently,
the crossed-module identity gives
$\partial(x)b\partial(x)^{-1}=b$ for every $b\in\Pi$.
Hyperbolic geometric curve groups are centre-free. The same argument
and conclusion are explicitly given in the proof of their Proposition
8.3, using Lemma 8.2 \cite{CF_HSS2026}. Hence $\partial=0$ and (C.3)
is exact. The same theorem for a base change identifies the geometric
kernel; the homomorphism to $E\times_G\pi_1(S')$ is the identity on
that kernel and on the quotient, and is an isomorphism.

For (C.4), an element in its kernel lies in $\Pi$ and centralizes
$\Pi$, so is trivial. Given $(a,g)$ choose $e\in E$ over $g$.
The automorphism $a\operatorname{conj}_e^{-1}$ is inner, so equals
$\operatorname{conj}_b$ for a unique $b\in\Pi$. Then $be$ maps to
$(a,g)$. This proves bijectivity, and compactness gives the topological
isomorphism. Thus the extension reconstructed in the margin is not
merely a larger object mapping onto $\pi_1(U)$ in this setting: the
displayed map is an isomorphism and supplies its inverse explicitly.

If $S$ is a $K(\pi,1)$, its $\pi_2$ is zero, so the same homotopy
sequence proves the centre-bearing cases. Over a trait the version
with section is SGA~1 XIII, Proposition 4.3 and Examples 4.4 (a section
exists \'etale locally and the claimed injection can be tested there).
An actual counterexample to the deleted unrestricted margin is
$S=\mathbb P^1_{\mathbb C}$ and
$U=\mathcal O_{\mathbb P^1}(-1)^\times$. The map
$(\ell,v)\mapsto v$ identifies $U$ with
$\mathbb A^2_{\mathbb C}\setminus\{0\}$; its complex analytification
retracts onto $S^3$ and has trivial fundamental group. Riemann existence
identifies its finite \'etale covers with finite topological covers.
Consequently $\pi_1(U)=1$, although its fibre $\mathbb G_m$ has group
$\widehat{\mathbb Z}(1)$. The boundary map from
$\pi_2(\mathbb P^1(\mathbb C))=\mathbb Z$ is multiplication by the
degree $-1$ of the line bundle and kills the fibre group. This computes
the obstruction which the margin had set aside.
\end{proof}

\begin{result}{C5}{Proved}{Unramified rational geometric sections extend}
\textup{[F11; \S3, manuscript pp.\ 32, 34--36, 41--42; scans 65,
70--74, 83--84.]}
Let $S$ be smooth integral of finite type over a characteristic-zero
field, with function field $K$, and let $U/S$ be a tower as in C1 for
which the exact geometric sequences of C4 hold. A rational section
$f_K$ extends to $S$ if and only if its induced lift
\[
 \lambda_f:G_K\longrightarrow\pi_1(U)
\]
kills $\ker(G_K\to\pi_1(S))$. Equivalently, it kills every divisorial
inertia subgroup of $G_K$ whose centre has codimension one in $S$.
\end{result}

\begin{proof}
The kernel in question is the closed normal subgroup generated by those
inertia groups. Indeed, a finite Galois extension of $K$ on which all
these inertia groups vanish has normalization \'etale in codimension
one over $S$; purity of the branch locus makes it \'etale everywhere.
Conversely an \'etale cover is unramified at each such valuation. This
identifies the finite quotients of the two proposed quotient groups,
and therefore identifies the profinite groups themselves. The purity
statement used here is the Zariski--Nagata theorem: a finite dominant
map with normal source and regular target, generically \'etale, has
branch locus either empty or of pure codimension one
\cite{CF_StacksPurity}.

If $\lambda_f$ kills the kernel, C2 applied at every height-one local
ring extends the section there. Finite presentation spreads each such
extension to a neighborhood of that point. The union with the original
domain of definition is an open set whose complement has codimension
at least two. C1 extends across that complement. The reverse implication
is the functorial factorization of fundamental groups associated with
the actual section. This proof specifies why the purity comparison
$\pi_1(S\setminus Z)\simeq\pi_1(S)$ by itself would not prove the
extension of a morphism: C1 supplies that geometric step.
\end{proof}

\begin{result}{C6}{Proved}{The normalizer extension and its sign}
\textup{[F14; \S3, manuscript pp.\ 39--41; scans 79--82.]}
Assume the hyperbolic case of C4 and let $q:S\to D$ be a boundary
section. Put $M=I_q/I_q^2$ and $N=M^\vee=q^*\mathcal O_{\bar U}(q)$.
For a chosen peripheral subgroup $L_q=\kappa_q(T)\subset\Pi$, the
normalizer $N_E(L_q)$ fits in an exact sequence
\[
 1\longrightarrow T\longrightarrow N_E(L_q)
       \longrightarrow G\longrightarrow1.                 \tag{C.5}
\]
With the positive Kummer inertia convention of C3 its extension class
is $c_1(N)$, and hence is $-c_1(M)$. More precisely, its image in
$H^2(G,\mu_n)$ corresponds under the comparison map to the Kummer
class of $N$ in $H^2(S,\mu_n)$; when $S$ is an \'etale $K(\pi,1)$ this
comparison is an isomorphism.
For type $(0,2)$ the same conclusion holds whenever the exact
sequence of C4 holds (in particular over the arithmetic $K(\pi,1)$
bases of C7): here $L_q=\Pi=T$, its arithmetic normalizer is all of
$E$, and the two distinguished boundary sections identify $U$
with the punctured normal line $N^\times$. Type $(1,0)$ has no cusp,
so no normalizer extension of this kind is asserted there.
\end{result}

\begin{proof}
The conjugacy class of $L_q$ is stable under $G$, because $q$ is defined
over $S$. For any $g\in G$, choose $e\in E$ over $g$ and $b\in\Pi$
with $eL_qe^{-1}=bL_qb^{-1}$. Then $b^{-1}e$ lies in the normalizer
and maps to $g$. Its kernel is $N_\Pi(L_q)=L_q$, by peripheral
self-normalization (Stix, Corollary 12). This proves surjectivity and
exactness, without asserting that a surjective homomorphism has a section.

The punctured formal neighborhood of $q$ has the same tame-cover
category as the punctured normal line $N^\times$. This statement means
the following concrete comparison. On an \'etale open where a local
equation $z_i$ for $q$ has been chosen, all tame covers are dominated by
the Kummer covers $z_i^{1/n}$; on overlaps their identifications are the
identifications of the corresponding roots of the normal transition
functions. SGA~1 XIII, Proposition 5.2 gives this local description;
descent of finite \'etale covers gives the gluing. It produces a
homomorphism $\pi_1(N^\times)\to E$ whose image normalizes $L_q$,
which is the identity on the base quotient and the peripheral injection
on the fibre. In particular the fibre map into $\pi_1(N^\times)$ is injective, since its composite with $E$ is $\kappa_q$. The homotopy sequence for $N^\times\to S$ is therefore short exact. The exact sequences consequently identify the tube
extension with (C.5). This is also the identification of Stix,
\S10.3, with the line $q^*\mathcal O(q)$ \cite{CF_Stix2008}.

Here is the sign computation with all conventions fixed. Choose frames
$e_i$ of $N$ and write $e_j=g_{ij}e_i$. For a vector
$v=e_iw_i=e_jw_j$, its scalar coordinates satisfy $w_i=g_{ij}w_j$.
Choose local roots $h_{ij}^n=g_{ij}$ and write $r_i^n=w_i$;
the transition of the root cover is $r_i=h_{ij}r_j$. On triple overlaps
the defect is $h_{ij}h_{jk}h_{ki}\in\mu_n$. This is exactly the
Kummer connecting cocycle $\delta_n([N])$. The deck transformation
$r_i\mapsto\zeta r_i$ is the positive inertia element specified in C3.
Thus the extension cocycle, computed by multiplying three local
trivializations, has sign $+\delta_n(N)$. The dual transition
functions are $g_{ij}^{-1}$, so $\delta_n(M)=-\delta_n(N)$.
Stix, \S6.4, makes the line-bundle extension/Chern comparison and its
sign explicit. This computation resolves the manuscript's alternatives
``conormal (ou normal, le diable sait)'' mathematically; it does not
claim to remove the author's written doubt by palaeography.

For example $N=\mathcal O_{\mathbb P^1}(-1)$ has degree $-1$ and
$N^\times=\mathbb A^2\setminus0$ as computed in C4. Its homotopy
boundary is the degree $-1$, agreeing with the displayed normal-line
convention. A dual convention for the generator of inertia would
reverse both signs, and must also reverse the Kummer character in C3.
\end{proof}

\begin{result}{C7}{Proved}{The arithmetic splitting and completed-unit criteria}
\textup{[F14--F15; \S3, manuscript pp.\ 40--42; scans 81--82, 85.]}
Let $S$ be smooth connected of finite type over a field finitely generated
over $\mathbb Q$, and suppose that $S$ is an \'etale $K(\pi,1)$.
Then $\operatorname{Pic}(S)$ is finitely generated, the integral continuous
comparison in degree two is an isomorphism
\[
 H^2(G,T)\simeq H^2(S,T),
\]
and Kummer theory gives
\[
 H^1(G,T)\simeq H^1(S,T)
       \simeq\widehat{\mathcal O(S)^*}
       :=\varprojlim_n\mathcal O(S)^*/\mathcal O(S)^{*n}.       \tag{C.6}
\]
For the boundary section of C6, (C.5) splits if and only if $M$,
equivalently $N$, is a trivial line bundle. A global basis of $M$ is
sufficient over any characteristic-zero base for which (C.5) exists.
Here a ``uniformizer along $q$'' means a basis of $I_q/I_q^2$; the
statement does not require a global equation of $q$ on an entire
neighborhood of its image.
\end{result}

\begin{proof}
Guralnick--Jaffe--Raskind--Wiegand, Proposition 6.1, states that the Picard
group of a normal scheme of finite type over a finitely generated field,
or over $\mathbb Z$, is finitely generated \cite{CF_GJRW1996}. Its
normality and arithmetic field hypotheses hold here. They cannot be
deleted, for example $\operatorname{Pic}^0$ of a positive-genus curve
over an algebraically closed field is divisible and usually infinite.

For every positive integer $n$, the Kummer sequence gives
\[
 0\to\mathcal O(S)^*/\mathcal O(S)^{*n}
  \to H^1(S,\mu_n)\to\operatorname{Pic}(S)[n]\to0,
\]
and an injection
$\operatorname{Pic}(S)/n\hookrightarrow H^2(S,\mu_n)$.
If $N_0$ annihilates the torsion of $\operatorname{Pic}(S)$, the
transition from an index divisible by $nN_0$ to $n$ kills the Picard
torsion term; the transition on units modulo powers is onto. Thus the
images in the system $H^1(S,\mu_n)$ eventually equal the unit subgroup,
and this inverse system is Mittag--Leffler. The system $H^0(S,\mu_n)$
is also Mittag--Leffler, since its terms are finite. Taking a factorial
cofinal sequence of indices, the derived-limit exact sequence for
continuous cohomology has zero $\varprojlim^1$ in the degrees at issue.
This is the sequence (0.2) of Jannsen \cite{CF_Jannsen1988}, with
the Mittag--Leffler hypotheses just verified. It follows that
$H^i(S,T)=\varprojlim_nH^i(S,\mu_n)$ for $i=1,2$ here.

For a $K(\pi,1)$ the finite local-system comparison
$H^i(G,\mu_n)\simeq H^i(S,\mu_n)$ holds in every degree by its
definition (equivalently by the contractibility of its universal
\'etale homotopy cover). Applying the same inverse-limit sequence on
both sides proves the asserted integral comparison. The inverse
limit of $\operatorname{Pic}(S)[n]$, with its multiplication transition
maps, is zero because its torsion has bounded exponent. This proves
(C.6). In particular we have justified the limit operation in scan 82,
including the potentially missing right exactness.

If the extension splits, its continuous class vanishes, hence all
$c_1(N)\bmod n$ vanish. The Kummer injections imply
$[N]\in n\operatorname{Pic}(S)$ for every $n$. In a finitely generated
abelian group the intersection of these subgroups is zero, so $N$ is
trivial. Conversely a nonzero normal vector along the entire section
defines its tangential fibre functor, or equivalently a section of
$N^\times\to S$ followed by the tube comparison in C6. It splits
(C.5). This construction does not use finite generation or the
$K(\pi,1)$ assumption.

Surjectivity of the normalizer map alone does not imply splitting.
For a concrete correction to scan 80, let $C/\mathbb Q$ be a smooth
proper genus-two curve and use $S=C$,
$\bar U=C\times C\to C$ and the diagonal as its single cusp. The
normal line is $T_C$, of degree $-2$. Its Kummer class modulo $3$ has
nonzero degree in $\mathbb Z/3$; hence (C.5) does not split, although
it surjects onto $G$. The curve $C:y^2=x^5-x$ supplies a specified
example (the polynomial has no repeated root, and the smooth proper
completion has genus two).
\end{proof}

\begin{result}{C8}{Proved}{The geometric torsor inside the completed torsor}
\textup{[F15; \S3, manuscript p.\ 42; scan 85.]}
Under the arithmetic hypotheses of C7, $\mathcal O(S)^*$ contains no
nonidentity element divisible by every positive integer in that group.
The completion map in (C.6) is injective. When $N$ is trivial, the set
of its nowhere-zero sections is a torsor under $\mathcal O(S)^*$.
Its tangential splittings inject into the set of $T$-conjugacy classes
of splittings of (C.5), a torsor under $H^1(G,T)$. Under these actions
the map of torsors is induced by the injective Kummer map
\[
 \mathcal O(S)^*\hookrightarrow\widehat{\mathcal O(S)^*}.
\]
For bases of the conormal line, rather than normal vectors, multiplying
the basis by $a$ changes this Kummer class by $-\kappa(a)$.
\end{result}

\begin{proof}
Let $F$ be the function field of $S$, finitely generated over $\mathbb Q$,
and let $k$ be the algebraic closure of $\mathbb Q$ in $F$; it is a
number field. Suppose $a\in F^*$ is an $n$th power in $F$ for every
$n>0$. On a normal projective geometrically integral model of $F/k$,
each prime-divisor valuation of $a$ is divisible by every positive
integer and hence is zero. Consequently $a$ and $a^{-1}$ are global
regular functions and $a\in k^*$. Every root of $a$ lying in $F$ is
algebraic over $k$, and hence lies in $k$. At each finite prime of
$k$, its valuation is again zero; thus $a\in\mathcal O_k^*$ and all
its roots in $k$ are units. Dirichlet's unit theorem (Milne, Theorem~5.1 \cite{CF_MilneANT}) makes this group
finitely generated; its intersection of all power subgroups is the
identity. Thus $a=1$. Applying this to
$\mathcal O(S)^*\subset F^*$ proves the first assertions. The same
argument covers transcendence degree zero by starting directly with
the number field.

For completeness, the classification of the splitting torsor is
explicit. Fix a splitting $s:G\to N_E(L_q)$. Every other splitting
has the unique form $s_c(g)=c(g)s(g)$ for a continuous crossed
homomorphism $c:G\to T$. The homomorphism identity for $s_c$ is
$c(gh)=c(g)+g c(h)$, and conjugation by $b\in T$ adds the coboundary
$b-gb$. Thus the classes form the $H^1(G,T)$-torsor asserted. If a
normal frame is multiplied by $a$, the corresponding root frames
differ by compatible roots of $a$; the cocycle is exactly
$g\mapsto g(a^{1/n})/a^{1/n}$. Therefore its class is $\kappa(a)$.
The kernel is the intersection of all power subgroups, already proved
trivial. Dualizing a conormal frame replaces $a$ by $a^{-1}$ and
changes the sign. It is the ordinary unit group, not its completion,
that acts on actual algebraic frames. This is the mathematical
correction supported by scan 85, where a possible circumflex remains
a palaeographic uncertainty.
\end{proof}

\begin{result}{C9}{Proved}{Divisorial failure of extension and tangential specialization}
\textup{[F16; \S3, manuscript pp.\ 46--47; scans 93--95.]}
Let $S$ be smooth integral in characteristic zero and let $f_K$ be a
rational section of a smooth pair as in C1. Let $S'$ be its maximal
open domain as a section of $U$. Every irreducible component of
$S\setminus S'$ has codimension one. At the generic point of such a
component $D_\alpha$, the proper extension meets a unique cusp, and
\[
 \mu_\alpha=\operatorname{ord}_{D_\alpha}(f_K^*z)>0
\]
is its intersection multiplicity, independent of a unit multiple of
the local cusp equation $z$. The sum
$\sum_\alpha\mu_\alpha D_\alpha$ is an effective Cartier divisor
because $S$ is regular.

On an open subset $W$ of the regular locus of $D_\alpha$ where the
compactified section and the cusp are defined, and where a uniformizer
$t$ of $D_\alpha$ has been fixed, write
$z\circ f=t^{\mu_\alpha}u$ with $u$ regular and invertible along $W$.
The resulting conormal isomorphism is
\[
 f_W^*(I_q/I_q^2)\xrightarrow{\sim}
            (I_{D_\alpha}/I_{D_\alpha}^2)^{\otimes\mu_\alpha},
       \qquad z\longmapsto u\,t^{\mu_\alpha}.               \tag{C.7}
\]
It produces tangential specialization of the geometric section at
every point of $W$, compatible with specialization inside $W$. Changing
$t$ to $at$ replaces the leading coefficient by
$a^{-\mu_\alpha}u$, and changes the associated Kummer class by
$-\mu_\alpha\kappa(a)$ in the scalar-coordinate convention.
\end{result}

\begin{proof}
If an irreducible component of the complement had codimension at least
two, remove the other components and work near its generic point. The
section is then defined outside a codimension-at-least-two closed set,
and C1 extends it, contradicting maximality. At a codimension-one
generic point, properness extends $f_K$ to $\bar U$. Its image cannot
lie in $U$, again by maximality and finite presentation. The finite
\'etale boundary has disjoint sections after an \'etale localization,
so exactly one is met, and C3 gives the stated multiplicity. The
divisorial sum is locally finite because the complement is a closed
subset of the noetherian scheme $S$; on a regular scheme Weil
divisors are Cartier. We do not identify it with the pullback of a
global divisor by an everywhere-defined compactified rational map:
such a map can still have a codimension-two indeterminacy in genus
zero.

Where $z\circ f=t^\mu u$ with $u$ a unit, reduction modulo the next
power of $I_D$ gives precisely (C.7). On the punctured normal lines
the corresponding map is, in scalar coordinates, $v\mapsto u v^\mu$.
It sends the source tangent vector with $t(v)=1$ to the target tangent
vector with $z(w)=u$. For every $n$, pulling the Kummer cover of the
target back by this map gives the $\mu$th-power Kummer cover multiplied
by the unit $u$. These compatible maps of finite covers define the
map of tangential fibre functors, hence the specialization of sections
of the exact fundamental-group extensions. Evaluation at a point of
$W$ commutes with these pullbacks of covers, proving the asserted
compatibility. The replacement $t'=at$ gives
$u'=a^{-\mu}u$ directly, and Kummer functoriality gives the final
formula. At a zero or pole of $u$, this construction no longer gives
a nonzero tangent vector; this is exactly the regular-and-invertible
restriction in the manuscript.
\end{proof}
